
\subsection{同余关系中的模数重缩放}
\label{subsec:modulo-rescaling}

设 $q,q'$ 为正整数模数，$r=q'/q$，并固定具有整数平移不变性的最近整数舍入规则 $R(x)=\lfloor x+1/2\rfloor$。对于整数 $h$，有 $R(x+h)=R(x)+h$。记 $\delta_x=R(rx)-rx$，则 $|\delta_x|\le1/2$。

\begin{table}[h!]
\centering
\begin{tabular}{|l|l|}
\hline
原关系 & 模 $q'$ 下的缩放关系 \\
\hline
$a\equiv b\pmod q$ & $R(ra)\equiv R(rb)\pmod{q'}$ \\
$a+c\equiv b+d\pmod q$ & $R(ra)+R(rc)\equiv R(rb)+R(rd)+\eta\pmod{q'}$ \\
$ac\equiv bd\pmod q$ & $R(rac)\equiv R(rbd)\pmod{q'}$ \\
\hline
\end{tabular}
\caption{同余关系的缩放，其中 $\eta=\delta_a+\delta_c-\delta_b-\delta_d$ 且 $|\eta|\le2$。}
\label{tab:rescaling}
\end{table}

\begin{proof}
若 $a=b+kq$，则 $ra=rb+kq'$，并且 $R(ra)=R(rb)+kq'$。因此，对同余式两侧的整体分别舍入，会精确保持同余，无论 $ra$ 是否为整数。相同论证也适用于整个乘积 $ac$ 和 $bd$。

对于分别舍入的加数，由 $a+c=b+d+kq$ 得到
\[
R(ra)+R(rc)-R(rb)-R(rd)=kq'+\delta_a+\delta_c-\delta_b-\delta_d.
\]
最后四项组成整数 $\eta$，满足 $|\eta|\le2$。当它们抵消时，特别是所有缩放后的加数均为整数时，关系精确成立；否则必须保留这个显式误差。
\end{proof}

这些结论是单向蕴含，一般不是等价关系。例如，取 $q=8$、$q'=2$，则 $a=0$ 与 $b=1$ 均舍入为零，但 $a\not\equiv b\pmod8$。舍入会丢失信息。此外，分别缩放乘积的两个因子会将整体尺度变为 $r^2$，不能替代将整个乘积只乘一次 $r$。

对于固定整数乘数 $s$，有用的恒等式为 $R(ra)s=ras+\delta_a s$，误差至多为 $|s|/2$。因此，在 LWE 相位 $b-\sum_i a_i s_i$ 中，分别舍入 $b$ 和各个 $a_i$ 会引入 $\delta_b-\sum_i\delta_{a_i}s_i$。后续模数切换正确性需要约束这个误差，而不能只使用没有误差定义的近似同余符号。

\subsubsection{示例}
\label{subsec:modulo-rescaling-ex}

假设有如下同余关系：

$b \equiv a\cdot s + \Delta \cdot m + e \bmod q$，\text{ }其中：$q = 30$，\text{ }$s = 5$，\text{ }$a = 10$，\text{ }$\Delta = 10$，\text{ }$m = 1$，\text{ }$e = 10$，\text{ }$b = 40$

$ $

首先，代入给定的示例值，验证上述同余关系是否成立：

$b \equiv a\cdot s + \Delta \cdot m + e \bmod 30$

$40 \equiv 10 \cdot 5 + 10 \cdot 1 + 10 \bmod 30$

$40 \equiv 70 \bmod 30$ 

$ $

该同余关系成立。

$ $

现在，希望将模数按 $30 \rightarrow 3$ 重缩放。根据 \autoref{tab:rescaling} 中介绍的重缩放原理，计算缩放后的值如下：

$q'= 3$, \text{ } $s = 5$, \text{ } $m = 1$

$\hat{a} = \Bigg\lceil a\cdot\dfrac{3}{30} \Bigg\rfloor = \Bigg\lceil 10\cdot\dfrac{3}{30} \Bigg\rfloor = 1$

$\hat{\Delta} = \Bigg\lceil \Delta\cdot\dfrac{3}{30} \Bigg\rfloor = \Bigg\lceil 10\cdot\dfrac{3}{30} \Bigg\rfloor = 1$



$\hat{e} = \Bigg\lceil e\cdot\dfrac{3}{30} \Bigg\rfloor = \Bigg\lceil 10\cdot\dfrac{3}{30} \Bigg\rfloor = 1$

$\hat{b} = \Bigg\lceil b\cdot\dfrac{3}{30} \Bigg\rfloor = \Bigg\lceil 40\cdot\dfrac{3}{30} \Bigg\rfloor = 4$

$ $

将模数按 $30 \rightarrow 3$ 重缩放后，得到的同余关系推导如下：

$\Bigg\lceil b\dfrac{3}{30} \Bigg\rfloor \equiv \Bigg\lceil s\cdot a \dfrac{3}{30} \Bigg\rfloor + \Bigg\lceil m \cdot \Delta \dfrac{3}{30} \Bigg\rfloor + \Bigg\lceil e \dfrac{3}{30} \Bigg\rfloor \bmod 3$

$\hat{b} \equiv \hat{a} \cdot s + \hat{\Delta} \cdot m + \hat{e}  \bmod 3$
 \text{ }（这是精确同余关系，因为所有缩放后的值均无小数部分。）

$4 \equiv 1 \cdot 5 + 1 \cdot 1 + 1 \bmod 3$

$4 \equiv 7 \bmod 3$

$ $

该例精确成立，因为 $ra$、$r\Delta$ 和 $re$ 都是整数，所有舍入误差均为零。若部分缩放后的项不是整数，误差抵消后仍可能精确成立；一般情况下，应保留 \autoref{tab:rescaling} 中相应的误差项。









